\section{Introduction}

Client-centric consistency describes what one application can observe across its own operations. A replica set can acknowledge a write at one member while a later read reaches another member with a different state. Read concern, write concern, read preference, and causal session state determine which histories the client can observe \cite{mongodb-read-isolation,mongodb-read-write-semantics}.

RQ1 asks how read concern, write concern, and causal sessions change RYW, MR, MW, and WFR under registered adversarial schedules. RQ2 separately asks how normal operation, secondary failure, primary election, and network partition change availability and latency. A recorded dependency or visibility rule is used to explain a counterexample. A stronger setting can reduce a returned counterexample by waiting or rejecting an operation; that outcome is reported separately from a successful history.

The planned contribution is an executable comparison rather than a claim about MongoDB in general. The checker consumes saved histories without a live database. The report separates a successful history that violates a predicate from an operation that was blocked, timed out, or failed before a successful history existed.
